Das Blei kühlt von $\vartheta_\text{Pb}$ auf die Endtemperatur $\vartheta$ ab. Die abgegebene Wärme erwärmt das Kalorimeter und das Petroleum von $\vartheta_1$ auf $\vartheta$ (Wärmeverluste vernachlässigt):
\begin{align}
 m_\text{Pb}c_\text{Pb}\left(\vartheta_\text{Pb}-\vartheta\right) &= \left(m_\text{Cu}c_\text{Cu}+m_\text{P}c_\text{P}\right)\left(\vartheta-\vartheta_1\right)
\end{align}
Mit $c_\text{Cu} = \cCuO$ und $c_\text{Pb} = \cPbO$ (FoTaBe S.~190) folgt
\begin{align}
 c_\text{P} &= \cPeF \\
   &= \frac{\mPb\cdot\cPb\cdot\left(\TPbO-\TmO\right)-\mCu\cdot\cCu\cdot\left(\TmO-\TaO\right)}{\mPe\cdot\left(\TmO-\TaO\right)} \\
   &= \cPe \approx \result{\cPeP{3}{2}}
\end{align}
Der Wert liegt im Bereich anderer organischer Flüssigkeiten (FoTaBe S.~191). Weil die Temperatur nur um wenige Zehntelgrad ansteigt, ist das Resultat empfindlich auf Messfehler.
